\documentclass{article}
\usepackage{graphicx} 
\usepackage[dutch]{babel}
\begin{document}
	\sffamily
	\begin{titlepage}
		\centering
		\vfill
		{\bfseries\Huge
			Verslag Tinlab Advanced Algorithms \\
			\vskip2cm
		}
		{\bfseries\Large
			T. Ravensbergen\\ 
			G. Bartes\\
			K. G. Razmjou\\
		}
		{
			\bfseries\normalsize
			69\\
			\vskip1cm
			\today\\
		}    
		\vfill
		\includegraphics[width=4cm]{logohr.png} % also works with logo.pdf
		\vfill
		\vfill
	\end{titlepage}
	\newpage
	\tableofcontents
	
	\newpage
	\section{Inleiding}
	In deze case study wordt voor de millitaire generalen van Duitsland iets gemaakt.
	
	\section{Requirements}
	
	\subsection{Requirements}
	Directe requirements van opdrachtgever:\\
	Na grondige analyse van het Nederlandse sluizenpark is gebleken dat renova-tie van een groot aantal sluizen noodzakelijk is.  Een eerste verkenning heeft onsgeleerd dat het gecombineerd renoveren en automatiseren van het Nederlandsesluizenpark een aanzienlijke verbetering kan opleveren t.a.v.:\\
	- veiligheid\\
	- efficientie\\
	- capaciteit\\
	- onderhoudskosten\\
	- duurzaamheid\\
	In het kader van het onlangs afgesloten klimaatakkoord heeft de Nederlandseoverheid  daarom  besloten  over te gaan tot een ingrijpende renovatie van dediverse sluizen die ons land rijk is. Op het ministerie van infrastructuur en waterstaat is helaas onvoldoende kennis van ict en systemen aanwezig om eenen ander uit te voeren. Wij vragen u een model (of een onderling samenhangend aantal modellen)aan  te leveren, opdat ontwerpen van verschillende, volledig geautomatiseerdesluizen in de toekomst gerealiseerd kunnen worden.\\\\
	Eigen inbreng van deze requirements:\\
	Wij gaan er van uit dat het volgende van ons verwacht wordt:\\
	Maak een model dat als template dient gebruikt te worden voor het automatiseren van verschillende soorten sluizen. Verder moeten overwegingen gemaakt worden die goed onderbouwd zijn.\\\\ Aangezien er van ons alleen een model verwacht wordt, zullen wij ons geheel focussen op de fundamentele werking van de sluis en hierbij zullen wij ons dus niet bezig  houden met fysieke eisen zoals veiligheidshekjes en borden. Onze focus ligt geheel op de werking van de sluis; elke state waar de sluis zich in mag bevinden en welke beslissingen de sluis moet maken op basis van bestaande protocols en benoemde eisen. \\\\
	Deze requirements zullen hieronder uitgewerkt worden, per sluisonderdeel, deze bestaande uit de sluisdeuren, de sloplichten, de waterpomp en de boten.\\
	
	\subsection{Sluisdeuren}
	De sluisdeuren.
	
	\subsection{Stoplichten}
	De stoplichten
	
	\subsection{Waterpomp}
	De waterpomp
	
	\subsection{Boten}
	De meeste sluizen die zich in Nederland bevinden zijn schutsluizen; deze sluizen zijn bedoeld om boten, zowel vrachtschepen als pleziervaart afhangend van de locatie van de sluis, te verwerken. Om deze reden gaan wij deze dus ook verwerken in ons model. Mocht een sluis niet bedoeld zijn om boten te verwerken, dan zou dit model alsnog toegepast kunnen worden opp desbetreffende sluis.
	
	Overige einsen op basis van eigen inbreng:\\
	
	
	
	\subsection{Specificaties}
	Vanuit deze requiremenst kunnen verdere specificaties opgesteld worden.
	
	\subsection{Het vier variabelen model}
	\subsubsection{Monitored variabelen}
	\subsubsection{Controlled variabelen}
	\subsubsection{Input variabelen}
	\subsubsection{Output variabelen}
	
	\subsection{Rampen}
	
	\subsubsection{Ramp 1}
	\begin{description}
		\item[Beschrijving]
		\item[Datum en plaats] 
		\item[Oorzaak]
		%Beschrijf wat er mis ging in termen van het vier variabelen model/requirements/specificaties
	\end{description}
	
	\subsubsection{Ramp 2}
	\begin{description}
		\item[Beschrijving]
		\item[Datum en plaats] 
		\item[Oorzaak]
		%Beschrijf wat er mis ging in termen van het vier variabelen model/requirements/specificaties
	\end{description}
	
	\subsubsection{Ramp 3}
	\begin{description}
		\item[Beschrijving]
		\item[Datum en plaats] 
		\item[Oorzaak]
		%Beschrijf wat er mis ging in termen van het vier variabelen model/requirements/specificaties
	\end{description}
	
	\subsubsection{Ramp 4}
	\subsubsection{Ramp 5}
	\subsubsection{Ramp 6}
	
	\section{Modellen}
	
	\subsection{De Kripke structuur}
	
	\subsection{Soorten modellen}
	
	\subsection{Tijd}
	
	\subsection{Guards en invarianten}
	
	\subsection{Deadlock}
	
	\subsection{Zeno gedrag}
	
	\section{Logica}
	
	\subsection{Propositielogica}
	
	\subsection{Predicatenlogica}
	
	\subsection{Kwantoren}
	
	\subsection{Dualiteiten}
	
	\section{Computation tree logic}
	
	\subsection{De computation tree}
	
	\subsection{Operator: AG}
	
	\subsection{Operator: EG}
	
	\subsection{Operator: AF}
	
	\subsection{Operator: EF}
	
	\subsection{Operator: AX}
	
	\subsection{Operator: EX}
	
	\subsection{Operator: p U q}
	
	\subsection{Operator: p R q}
	
	\subsection{Fairness}
	
	\subsection{Liveness}
	
	\newpage
	
	\newpage
	\bibliography{references}
	\bibliographystyle{plain}
\end{document}


